\section{Phương pháp đánh giá}\label{sec:eval-method}

\subsection{Hai nguồn dữ liệu}
Nguồn thứ nhất là nhật ký các buổi học thử nghiệm thật, đọc từ bản sao lưu cơ sở dữ liệu ngày 16/11/2026 bằng một lệnh chỉ đọc. Bản sao gồm $43$ buổi của năm tài khoản học sinh thử nghiệm với $14\,078$ sự kiện; mỗi sự kiện ghi một việc đã xảy ra (học sinh nhập phản hồi, gia sư đưa gợi ý, bộ chấm trả kết quả) và chỉ được bổ sung, không bị sửa. Tài khoản Vuong chiếm $79{,}7\%$ số sự kiện và $92{,}7\%$ số câu hỏi đã hiển thị ($523$ trên $564$ câu) vì học nhiều nhất, nên các mục kết quả lấy tài khoản này làm nguồn chính: $12$ buổi có ít nhất một câu được giải quyết, gồm sáu buổi Toán chương Số thập phân, bốn buổi Toán chương Diện tích và hai buổi Tiếng Anh Unit~11. Bốn tài khoản còn lại chỉ góp vào các phép đếm về mô hình ngôn ngữ và bộ kiểm chứng ở mục~\ref{sec:eval-model}, nơi kết quả không phụ thuộc vào việc ai học. Các số liệu được tính lại bằng chương trình đọc nhật ký và chỉ xuất số liệu tổng hợp, không xuất phản hồi học sinh hay tên tài khoản.

Nguồn thứ hai là các bộ câu hỏi và tình huống soạn sẵn, được lưu lại bằng thẻ git trước khi sửa phần hệ thống mà chúng đo, gọi là \emph{đóng băng}. Bộ đóng băng cho phép so sánh công bằng giữa các cách làm. Bảng~\ref{tab:eval-sets} liệt kê các bộ, do ai viết, dùng để làm gì và có độc lập với thiết kế không; điều này quyết định bộ nào được dùng làm bằng chứng kết luận.

\begin{table}[H]
\centering
\caption{Các bộ dữ liệu đánh giá. Bộ độc lập là bộ chưa từng được xem để sửa thiết kế. Nhãn truy xuất là mọi đoạn của bài trả lời câu hỏi, do người phát triển và một trợ lý AI soạn, chưa có chuyên gia thẩm định.}\label{tab:eval-sets}
\small
\begin{tabular}{>{\raggedright\arraybackslash}p{2.7cm}r>{\raggedright\arraybackslash}p{5.2cm}>{\raggedright\arraybackslash}p{4.3cm}}
\toprule
Bộ & Cỡ mẫu & Nội dung và cách dựng & Độc lập với thiết kế? \\
\midrule
Bộ chỉnh & $30$ (trong bài $18$) & Câu hỏi truy xuất dùng để chỉnh ngưỡng định tuyến & Không \\
Bộ 1 & $23$ ($14$) & Viết sau bộ chỉnh; các câu trượt đã được đọc khi thiết kế định tuyến & Không \\
Bộ 2 & $30$ ($20$) & Viết sau khi thiết kế xong, nhiều câu hỏi bài trước hơn ($12$ so với $7$) vì định tuyến nhắm vào loại câu này; đóng băng trước khi sửa mã & Có \\
Bộ 3 & $43$ ($30$) & Câu viết như trẻ gõ không dấu, viết tắt; dùng để sửa ở lần cải tiến cuối & Không \\
Bộ 4 & $37$ ($22$) & Câu viết như trẻ gõ; đóng băng trước lần cải tiến cuối & Có \\
$13$ vai học sinh khó & $31$ lượt & Kịch bản hội thoại nhóm soạn, một bài tập & Có \\
Bộ chấm 1 và 2 & $36$ và $35$ câu & Câu trả lời điển hình của trẻ, gán nhãn đúng sai bằng tay & Bộ 2 độc lập; Bộ 1 đã xem khi sửa \\
Chuỗi sự kiện & $2\times500$ & Chuỗi ngẫu nhiên kiểm công thức mức hiểu & Có \\
\bottomrule
\end{tabular}
\end{table}

Vì Bộ chỉnh, Bộ 1 và Bộ 3 không độc lập, kết luận về truy xuất dựa chủ yếu vào Bộ 2 và Bộ 4; các bộ kia chỉ cho thấy chiều của tác dụng. Mô hình nhúng \texttt{gemini-embedding-001} đo độ gần về nghĩa, khoảng tin cậy tính bằng lấy mẫu lại $10\,000$ lần, kiểm định McNemar so hai cách chạy trên cùng một bộ câu, và khoảng Wilson $95\%$ cho các tỉ lệ nhỏ.

\subsection{Các chỉ số dùng trong chương}
Hit@3 là tỉ lệ câu hỏi mà đoạn đúng nằm trong ba kết quả đầu; MRR là trung bình của $1/\text{thứ hạng}$ của đoạn đúng đầu tiên (luôn đứng đầu thì bằng $1$, không có thì bằng $0$). Hit@3 và MRR chỉ tính trên các câu có đoạn đúng trong kho, tức câu về bài đang học và bài trước. Độ gần cosine đo hai đoạn văn giống nhau về nghĩa trong khoảng từ $0$ đến $1$. Khoảng tin cậy $95\%$ là khoảng mà giá trị thật nhiều khả năng nằm trong. $p$ của kiểm định McNemar cho biết chênh lệch giữa hai cách chạy có thể chỉ do ngẫu nhiên hay không; $p<0{,}05$ nghĩa là chênh lệch khó là ngẫu nhiên. Điểm hiểu bài $M$ từ $0$ đến $100$ ước lượng học sinh hiểu một bài đến đâu, kèm mức hiểu từ $0$ đến $5$. Hối tiếc tích lũy của bộ chọn nước đi là tổng phần thưởng bị mất so với luôn chọn nước tốt nhất. Một câu có bốn cách kết thúc: \emph{tự làm} (đúng không cần gợi ý), \emph{có gợi ý} (đúng sau một đến năm bậc), \emph{làm mẫu} (phải đến bậc sáu hoặc bảy) và \emph{chưa giải quyết}.

\subsection{Phạm vi của các kết luận}
Các số trong chương đo cách hệ thống hoạt động, không đo hiệu quả học tập của trẻ em. Các buổi do nhóm phát triển thực hiện bằng tài khoản học sinh thử nghiệm Vuong, chưa có trẻ em tham gia, và các buổi Toán tháng 10 chạy trên những bản gia sư cũ hơn bản đang triển khai. Chương vì thế chỉ kết luận về hai điều: hệ thống có giữ các quy tắc đã đặt không, và nó khác chatbot thông thường ra sao trong cùng tình huống.
